\documentclass[11pt,a4paper]{article}

\usepackage[margin=2.3cm]{geometry}
\usepackage[T1]{fontenc}
\usepackage[utf8]{inputenc}
\usepackage{lmodern}
\usepackage{setspace}
\usepackage{hyperref}
\usepackage{booktabs}
\usepackage{enumitem}

\title{Interpretable Neuro-Functional Music AI:\\
\large A Position Paper on the \textit{lif\_lova} LIF Architecture}
\author{Draft for Workshop Submission}
\date{\today}

\begin{document}
\maketitle
\onehalfspacing

\begin{abstract}
This position paper argues that the current \textit{lif\_lova} architecture represents a productive
alternative to mainstream black-box music generation. Instead of treating composition as token
prediction, it combines interpretable neural dynamics (leaky integrate-and-fire network states)
with explicit music-theoretical constraints (functional harmony, modality, register, and style-conditioned
voicing). We frame this as an \emph{interpretable neuro-functional paradigm}: post-symbolic in state formation,
symbolic in harmonic decision, and performative in real-time control. We outline why this architecture is
relevant to contemporary musicology and music AI, identify current limitations, and propose a concrete
research agenda with measurable evaluation criteria for interactive musical systems.
\end{abstract}

\section{Thesis}

The central claim is that \textit{lif\_lova} should be understood as an \emph{instrumental AI system},
not merely a generative model. Its value lies in three properties that are often separated in current research:
\begin{enumerate}[label=\arabic*.]
    \item \textbf{Causal transparency}: the pathway from signal to note is inspectable.
    \item \textbf{Music-theoretical intelligibility}: harmonic behavior is articulated through explicit functional and modal logic.
    \item \textbf{Performative responsiveness}: system behavior is directly controllable in live contexts.
\end{enumerate}

This makes the architecture especially relevant to computational musicology, live electronic performance,
and practice-based AI research.

\section{Context in Music AI Research}

Recent foundation-model approaches to music generation have improved stylistic plausibility and long-context
prediction. However, they commonly face recurring limitations in artistic research settings:
\begin{itemize}
    \item low interpretability of generation decisions,
    \item weak causal controllability during performance,
    \item heavy dependence on large copyrighted corpora,
    \item and unclear mapping between model internals and established music-theoretical constructs.
\end{itemize}

By contrast, \textit{lif\_lova} uses compact, designed dynamics rather than large corpus pretraining.
The system offers explicit intermediate representations (activity bins, harmonic function assignment,
modal pitch filtering, style-conditioned voicing), making analytical and pedagogical use feasible.

\section{Musicological Relevance}

\subsection{Function as Process}

In many symbolic systems, tonic/subdominant/dominant labels are imposed externally as static templates.
In \textit{lif\_lova}, harmonic function is inferred from evolving neural energy distributions and local contour.
This aligns with process-based musicological readings where function is enacted over time,
not merely declared by chord symbols.

\subsection{Modal Grammar with Dynamic Agency}

The architecture constrains pitch material by per-scene modal selection (Ionian through Locrian),
while timing, density, and articulation remain state-dependent. This realizes a productive split between:
\begin{itemize}
    \item \emph{grammar} (mode, key center, note-range framing), and
    \item \emph{agency} (network-driven activation and release).
\end{itemize}

\subsection{Topology as Form-Generating Constraint}

Network topology (ring, fully connected, feedforward, sparse random, small-world) acts as a latent formal parameter.
Each topology shapes recurrence, burst behavior, and continuity, thereby influencing perceived phrase and texture.
From a musicological perspective, topology can be treated as a model of form-bearing tendency.

\section{Why This Matters Beyond One System}

The architecture models a broader design pattern:
\begin{center}
\textbf{continuous dynamics} $\rightarrow$ \textbf{interpretable musical state} $\rightarrow$ \textbf{performable output}
\end{center}

This pattern may be preferable for creative AI tools where authorial control, accountability,
and communicable analysis are as important as novelty.

\section{Current Limitations}

Despite its strengths, the present architecture has identifiable research gaps:
\begin{enumerate}[label=\arabic*.]
    \item \textbf{Limited long-horizon memory}: behavior is strong at short and meso timescales but less explicit at phrase-scale planning.
    \item \textbf{Crisp functional boundaries}: hard or near-hard functional transitions may reduce nuanced ambiguity.
    \item \textbf{Discrete style states}: style modes are categorical rather than continuous manifolds.
    \item \textbf{Evaluation gap}: current success criteria are mostly experiential and engineering-centric.
\end{enumerate}

\section{Research Agenda}

\subsection{A. Multi-timescale Formal Memory}

Introduce slow control states (for example cadence pressure, repetition debt,
tension trajectory) that evolve over tens of seconds or minutes.
This would support longer formal arcs without abandoning local responsiveness.

\subsection{B. Probabilistic Functional Inference}

Replace single-label function assignment with a distribution over T/S/D classes,
then smooth in time with confidence-aware hysteresis.
This enables analytically meaningful functional ambiguity and gradual reinterpretation.

\subsection{C. Continuous Style Manifolds}

Generalize style presets into interpolable voicing spaces.
Instead of switching Pop/Rock/Jazz/Blues as discrete categories,
performers could navigate trajectories between them.

\subsection{D. Co-Adaptive Session Learning}

Add a lightweight online adaptation layer to tune density, register preference,
and articulation thresholds based on performer interaction,
while keeping the harmonic rule core transparent.

\section{Evaluation Protocol for Interactive Music AI}

We propose a mixed-method protocol that combines quantitative telemetry and qualitative judgment.

\subsection{Metrics}

\begin{table}[h]
\centering
\begin{tabular}{@{}ll@{}}
\toprule
Dimension & Example operationalization \\
\midrule
Controllability & Parameter perturbation causes predictable output shifts \\
Harmonic intelligibility & Expert ratings of functional clarity and modal consistency \\
Temporal coherence & Segment-level continuity and phrase-shape stability \\
Expressive range & Coverage of density, register, and articulation states \\
Repeatability vs novelty & Controlled reruns with bounded divergence \\
Performer fit & Time-to-intent and correction burden in live sessions \\
\bottomrule
\end{tabular}
\caption{Core dimensions for evaluating instrumental AI systems.}
\end{table}

\subsection{Study Design}

\begin{enumerate}[label=\arabic*.]
    \item Recruit expert users (improvisers/composers) and technically proficient users.
    \item Compare three conditions: fixed-rule baseline, current \textit{lif\_lova}, and enhanced variant.
    \item Collect synchronized logs (control changes, topology, note events, audio descriptors).
    \item Run post-session annotation interviews with timeline replay.
    \item Triangulate metrics with expert panel listening.
\end{enumerate}

\section{Position Statement}

The field should treat interpretable, performable hybrid systems as first-class music AI research objects,
not as transitional artifacts between symbolic systems and large generative models.
\textit{lif\_lova} demonstrates that explicit musical knowledge and neural dynamics can be co-compositional,
yielding systems that are analyzable, steerable, and artistically credible.

\section{Conclusion}

The current \textit{lif\_lova} LIF architecture is already a compelling contribution to modern music AI discourse:
it embodies an alternative research trajectory centered on transparency, theory-informed behavior,
and live musical agency. The next step is not to replace this architecture with opaque end-to-end generation,
but to deepen it through multi-timescale memory, probabilistic harmonic reasoning,
continuous stylistic control, and rigorous interactive evaluation.

\section*{Suggested Citation Placeholder}

Author(s). \emph{Interpretable Neuro-Functional Music AI: A Position Paper on the lif\_lova LIF Architecture}.
Workshop draft, 2026.

\end{document}